\begin{abstract}
\noindent
Can countries meet the basic social needs of their populations while staying within a fair per-person share of the climate boundary? This paper assembles a country-year panel of \nNCountries{} countries for \nFirstYear--\nLastYear{} and compares four social thresholds (life expectancy, under-5 mortality, access to electricity and employment) with consumption-based \COtwo{} emissions, measured against an equal-per-capita boundary of \nBoundary~\tpc. None of the \nNCons{} countries with consumption-based emissions estimates met all four thresholds while emitting within the boundary in any year. All \nNCrossers{} of them that crossed the social foundation during the period did so while overshooting, as did all \nNCrossersSust{} when crossing requires the foundation to be met for three consecutive years. With production-based emissions, which are available for all \nNCountries{} countries in the panel, only \nProdSafeCountriesAll{} were in the safe and just space, in \nProdSafeYearsAll{} country-years in total, and in those years both exceeded the boundary once emissions embodied in imports are counted. I summarise the cross-country relationship with a descriptive ``carbon requirement'': the \COtwo{} per person at which half of countries meet the social foundation. It is not an estimate of what any single country needs. It fell from about \nReqStart~t in \nShareYearStart{} to \nReqEnd~t in \nLastYear{} (95\% CI \nReqEndLo--\nReqEndHi), a decline of \nReqTrend\% a year, yet remained \nReqEndMultiple{} times the boundary; for the health thresholds alone it fell to \nReqHealthEnd~t. Its level depends on how the social foundation is defined: without the employment threshold it was \nRobReqNoEmp~t, and with secondary enrolment added \nRobReqEdu~t. Its decline, and its position above the boundary, held in every specification. Absolute decoupling of consumption-based emissions from income spread from \nDecAbsEarly{} countries in 2001--2010 to \nDecAbsLate{} of \nDecN{} in 2011--2022. Yet only \nDecAboveOnTrack{} of the \nDecAbsAbove{} decouplers above the boundary cut emissions fast enough to reach it by 2050: the median cut was \nDecMedianAchieved\% a year against \nDecMedianRequired\% required. Even at the fastest rate any of them achieved, only \nScenFastest{} would arrive in time. The four configurations of shortfall and overshoot were persistent, and no country moved into the safe and just space. Development has become much less carbon-intensive, but not nearly fast enough to bring the carbon cost of meeting these basic needs within the climate boundary. The analysis covers one planetary boundary and four social indicators, and its estimates are descriptive associations, not causal effects.

\medskip
\noindent\textbf{Keywords:} planetary boundaries; social foundation; doughnut economics; consumption-based emissions; decoupling; carbon intensity of well-being; sustainable development

\medskip
\noindent\textbf{JEL:} O15, Q01, Q56, I15
\end{abstract}
